%% ================================================================================================
%%  1. Introduction
%%  Notes: research-agenda.md §0 (question, hypotheses, requirements), §1 (answer in brief)
%% ================================================================================================
\section{Introduction}\label{sec:intro}

\subsection{The question}

A group has to choose one option among several: a site, a design, a budget allocation, a policy. The
group wants a \emph{consensus} choice, one that everyone can live with, and it wants a procedure that
produces one. This note asks what the research literature says about the optimal way to vote in that
situation, and turns the answer into a procedure that can be implemented. We work under three
hypotheses, set by the use case:
\begin{description}
  \item[H1 (open agenda)] any voter may propose as many options as they wish;
  \item[H2 (full ballots)] every voter expresses a view on \emph{all} the options proposed;
  \item[H3 (veto rights)] voters have a right of veto.
\end{description}
The procedure must meet two requirements:
\begin{description}
  \item[R1 (best compromise)] the option chosen is the best compromise, in a sense to be made precise;
  \item[R2 (no deadlock)] the community never gets stuck, so the veto must not be too stringent.
\end{description}

The two requirements pull against each other. The traditional consensus rule, unanimity, gives every
member a veto: nothing passes to which anyone objects. It protects every minority, which is the spirit
of R1, but a single objector can block every change, and in large or divided groups someone always
objects (Section~\ref{sec:limits}). The Polish--Lithuanian \emph{liberum veto} is the textbook
case of a constitution paralysed by individual vetoes \citep{Rohac2008}. Majority rule, at the other
extreme, always decides (R2) but can impose on a large minority an outcome it finds unacceptable, which
is not a compromise. The question is therefore one of design: how much veto power, held by whom, over
what, and how to aggregate the rest of the ballot.

\subsection{The answer in brief}

\begin{enumerate}
  \item \textbf{No rule is optimal in every respect.} Every rule that is not dictatorial can be
  manipulated \citep{Gibbard1973,Satterthwaite1975}, majorities can cycle \citep{Condorcet1785}, and
  the majority principle cannot be reconciled with an effective minority veto: some profiles force a
  choice between the option a majority prefers and the one a large minority can accept
  (Proposition~\ref{prop:majority-veto}). A procedure is a choice of trade-off, which should be explicit
  and tunable.
  \item \textbf{The literature does not support hard unanimity.} Theory, experiments and institutional
  studies agree that unanimity raises delay, strengthens the status quo and gives holdouts bargaining
  power. Veto rights work better when they are separated from the approval of options and limited to
  blocking outcomes that are unacceptable
  \citep{BoutonLlorenteSaguerMalherbe2018,MiettinenVanberg2025,Tsebelis2002,DziudaLoeper2016}. In a large
  assembly, unanimity with the status quo as default deadlocks almost surely
  (Proposition~\ref{prop:gridlock}; 99.6\% of our simulated elections with 101 voters).
  \item \textbf{An open agenda requires a rule that is independent of clones.} When anyone can add
  options, a faction can flood the agenda with variants of its favourite. Borda's rule, fallback
  bargaining and the proportional veto core reward such flooding; the Schulze, Ranked Pairs and Split
  Cycle methods do not \citep{Tideman1987,Schulze2011,HollidayPacuit2023splitcycle}.
  \item \textbf{The procedure we recommend, \QVC($\theta$)}, Condorcet with a blocking quota, uses a
  single ballot on which each voter ranks all the options and marks those they find unacceptable.
  Options vetoed by at least a share $\theta$ of the voters are excluded; if that excludes them all, the
  least-vetoed are kept; the Schulze winner among the options left is elected. By default
  $\theta = 1/3$. The procedure always decides in one round, the status quo has no default advantage, the
  winner is acceptable to more than two thirds of the voters whenever some option is, the option that
  beats every other acceptable option in pairwise majority wins, and adding variants of an option changes
  nothing (Section~\ref{sec:procedure}).
  \item \textbf{The price is strategic vetoing.} A blocking minority can block strategically as well as
  sincerely. In our simulations, the share of elections that a coalition of 10\% of the voters can
  manipulate rises from 5--7\% under Schulze to 11--27\% under \QVC(1/3) in structured electorates
  (Section~\ref{sec:simulation}). Named and reasoned vetoes, a quota no lower than $1/4$ and at most one
  revision round contain this risk (Section~\ref{sec:safeguards}).
\end{enumerate}

\subsection{What this note adds and how it is organised}

The note starts from a literature review of optimal voting with veto rights produced with the
Consensus search engine in October 2026 (summarised in the repository's \texttt{literature/} folder). That
review covers committee bargaining and the comparison of unanimity with majority-with-veto. We add the
social-choice literature on Condorcet methods, clones and proportional veto, the evidence from
institutions that already vote this way, formal results on deadlock and on the procedure, and a Monte
Carlo study of 21 rules on 2,000 simulated elections per setting. Every number in the note comes from
code in the repository (\texttt{simulations/python/}), with checks against known results.

Section~\ref{sec:setting} sets out the model and the definitions. Section~\ref{sec:limits} states what
no procedure can achieve. Section~\ref{sec:veto} reviews the research on veto rights,
Section~\ref{sec:aggregation} the rules that aggregate full ballots. Section~\ref{sec:procedure} defines
and analyses the procedure, Section~\ref{sec:simulation} compares it with the alternatives by
simulation, and Section~\ref{sec:discussion} discusses limits, open questions and implementation.
Proofs are in Appendix~\ref{app:proofs}, algorithms in Appendix~\ref{app:algorithms} and the details of
the simulation in Appendix~\ref{app:simulation}.
